\section{A free partial source with no admissible retained completion}
\label{sec:136-free-partial-noncompletion}

\paragraph{Archive and exact scope.}
The conditions are cote136, physical pp.~117 and119,
footers116 and118, scans117 and119; the common-source question
is physical p.~121/footer120/scan121.
This applies the already proved construction
ARCH-136-DEGREE-TWO-FREE relatively under the single fixed
source of ARCH-136-LOW-TATE-RETAINED-COMPLETION.
It exhibits a bad choice of partial source. It does not show that
every partial source is bad, exclude a different common source,
or resolve the original existential forward-construction question.

\paragraph{ARCH-136-FREE-PARTIAL-NONCOMPLETION.}
There are an admissible integral theory \(C\), a strict source
\(H\), and two strict maps \(h_0,h_1:H\rightrightarrows C\)
such that every group \(H_n^{w,q}\) is free over \(\mathbb Z\),
\(H_0=S\), structural weight zero is \(\mathbb Z[0]\), and
\[
\begin{gathered}
 H^0(H_n^{w,*})=S^w,\qquad H^1(H_n^{w,*})=0
                         \quad\text{in every weight},\\
 H^q(H_n^{w,*})=0
                  \quad(q\geq3,\ w\geq\lambda_q),\\
 H_i(NH_*^{w,q})=0
                  \quad(q\geq0,\ i\geq0,\ w\geq\mu_q).
\end{gathered}                                         \tag{136.FPN1}
\]
The bounds are independent of simplex dimension and horizontal
degree. The remaining second cohomology is integer torsion.
There is no factorization
\[
 H\longrightarrow E\rightrightarrows C
                                                        \tag{136.FPN2}
\]
retaining \(h_0,h_1\) with \(E\) admissible.
Consequently \(H\) cannot have the missing uniform second-cohomology
tail. Both targets are nevertheless the same admissible theory.

\begin{proof}
Let \(A_{\leq2}\rightrightarrows C_{\leq2}\) be the original
single block of LOW-TATE-RETAINED-COMPLETION, and let \(A\)
be its full left Kan extension. That result, including its
coefficient shift (136.LTQ26)--(136.LTQ28), proves that the
two maps of this fixed \(A\) have no admissible retained
intermediate, with no bound left on that intermediate.
We check explicitly that \(A\) satisfies the initial hypotheses
needed for the free degree-two concentration construction.

Write \(D\) for the six-generator algebra (136.LTQ4).
Its new generators have zero boundary on the relative two-cell.
The free presentation therefore gives, in every simplex,
\[
 A_n\cong L(F)_n\otimes_S
          \bigotimes_{\alpha:[n]\twoheadrightarrow[2]}D_\alpha.
                                                        \tag{136.FPN3}
\]
The index consists of order-preserving surjections; it is finite
for each \(n\). A face that makes the composite non-surjective
sends the new generators to zero, and all other faces and
degeneracies send them to the corresponding copy. This verifies
the presentation and its left Kan universal property. It introduces
no relations between distinct copies beyond the common coefficients.

Each \(L(F)_n\) is a free tensor of even zero-form unit disks.
Each \(D_\alpha\) is a free graded PD algebra over \(S\).
Thus every group in (136.FPN3) is free abelian. Rationally
the triangular coordinates of the LTQ proof turn each \(D_\alpha\)
into free unit disks, so \(A\) has rational coefficient cohomology
in every simplex and weight. Its degree-zero forms are precisely
those of \(L(F)_n\), giving \(H^0=S\). A possible new one-cycle
has the form
\[
                  a+\sum_\alpha f_\alpha y_\alpha,
       \qquad a\in L(F)_n^1,\quad f_\alpha\in L(F)_n^0.
\]
The independent \(x_\alpha\) coefficients in its differential
are \(p f_\alpha\). The old zero-form module is integer-torsion-free,
so every \(f_\alpha=0\). Its one-cycles and one-boundaries are
therefore those of \(L(F)_n\), proving \(H^1=0\) in every
weight. Simplex zero and structural weight zero have the required
literal coefficient values.

Now run DEGREE-TWO-FREE under this \(A\), retaining its two
maps to \(C\). The targets already have simplex zero \(S\),
so its preliminary target replacement is unnecessary.
The proof of that construction uses its initial constant source
only to establish the following induction hypotheses: freeness of
the form groups, simplex zero \(S\), structural weight zero,
rational coefficient cohomology, and \(H^0=S,H^1=0\) in
every weight. These have just been verified for \(A\).
Every later operation is a free pushout under the existing source.

More explicitly, choose the same fixed target-dependent horizontal
cutoffs and attach the horizontal unit disks. At a fixed simplex
their boundary inclusions split on the free basis, so the resulting
graded extension is free. The unit-disk block calculation preserves
the stated low cohomology. For each \(q\geq3\), use the same
fixed vertical cutoff, and attach the reduced-representable Tate
pairs that kill all old cycles in that range. Rational comparison
makes those classes integer torsion. Target vertical exactness
supplies their primitive pairs; the two new form degrees
\(q-1,q-2\) are positive, so their vertex restrictions are zero.
This is the actual relative representable pushout under the old
source, and both target maps extend. No condition on its old
first skeleton is imposed.

For \(q\geq4\) this does not change degrees zero and one.
For \(q=3\), the independent new \(c_\alpha\) coefficient
of a possible one-cycle is \(N_\alpha f_\alpha\); freeness
of the old zero-form group forces \(f_\alpha=0\).
Thus no RH1 quotient is needed. The monomial complements are
free abelian and split off the old graded groups, so the compatible
union of bases remains a basis at every limit stage.
All these are precisely the preserved steps proved in
DEGREE-TWO-FREE; none requires that the starting object was
the constant coefficient source.

Alternate the horizontal and \(q\geq3\) vertical steps with those
fixed cutoffs. The same finite-support class-killing argument gives
the strict limit source \(H\), a map \(A\to H\), and the two
retained maps to \(C\), with (136.FPN1). Rational coefficient
cohomology persists, so its second cohomology is integer torsion.
The coefficient identifications are the specified ones; all maps
preserve them. The inherited construction has no remaining
admissibility condition except the uniform degree-two vertical tail.

If (136.FPN2) existed, its composite with \(A\to H\) would
give the admissible retained intermediate excluded by
LOW-TATE-RETAINED-COMPLETION. This proves noncompletion.
If \(H\) itself had the missing second-cohomology tail, it
would be admissible and would serve as that intermediate.
Therefore that tail fails for this \(H\).
This argument concerns these retained maps and this partial source;
the identical target \(C\) itself is still a common source for
the target pair.
\end{proof}

\paragraph{Verification status.}
The author reread the complete DEGREE-TWO-FREE proof and checked
each initial and preserved hypothesis, the pointwise left Kan
presentation, the two low-cohomology identities, every target
extension and cutoff, and the final factorization contradiction.
Reviewers cote15 and cote1348 each completed a whole saved-proof
reading, including fresh complete readings of DEGREE-TWO and
DEGREE-TWO-FREE, and reported PASS. The root also completed a whole saved-proof reading and reported
PASS, with contributor scope for proposing this relative-constructor
corollary. All actual checks are recorded separately.
No conclusion about the original existential main gate is inferred.
